% m-conclusiones.tex
% !TEX root = z-main.tex


\section{Conclusiones}

El desarrollo de modelos tridimensionales del sitio arqueológico Kaminaljuyú permitió alcanzar los objetivos planteados y responder a la necesidad de contar con herramientas digitales que favorezcan la preservación y difusión del patrimonio arqueológico guatemalteco. La aplicación del modelado y la animación \gls{3d} demostró ser una estrategia efectiva para representar digitalmente estructuras que, en la actualidad, se encuentran ocultas o deterioradas, contribuyendo así a su conservación visual y educativa.

La fidelidad geométrica y la coherencia visual de los seis montículos modelados confirmaron la viabilidad técnica del proceso metodológico implementado. El uso del software Blender, junto con la validación arqueológica y la gestión digital colaborativa, garantizó resultados consistentes y replicables. La integración de texturas realistas mejoró la percepción visual de las estructuras, facilitando su comprensión por parte de los usuarios.

La evaluación de los modelos en entornos de realidad aumentada evidenció su efectividad como medio de exploración e interpretación arqueológica. Los usuarios y especialistas validaron su utilidad para visualizar el sitio en contexto y reconocieron el potencial educativo de la herramienta.

Los resultados alcanzados demuestran que la combinación de modelado tridimensional y realidad aumentada constituye un recurso tecnológico aplicable a la investigación, conservación y divulgación del patrimonio. El trabajo desarrollado sienta las bases para futuras implementaciones en otros sitios arqueológicos del país, promoviendo la digitalización y accesibilidad del legado histórico guatemalteco mediante herramientas tecnológicas contemporáneas.

